% Appendix - ICML style

\section{Dataset Details}
\label{app:data}

\paragraph{WorldTrace.}
WorldTrace is a large-scale, open-source trajectory dataset derived from OpenStreetMap GPS traces.
We use trip segments from Detroit and Columbus, filtered for minimum length and valid map-matching.
Each segment retains the original Unix epoch departure time, enabling temporal conditioning.

\paragraph{Road graph construction.}
We extract the drivable road network from OpenStreetMap using OSMnx.
Nodes are projected to a $1024 \times 1024$ grid coordinate system, and edges retain their road type classification (motorway, primary, secondary, residential, service).
We aggregate road types into three tiers: major (motorway, primary), minor (secondary, tertiary), and service (residential, service).

\paragraph{Map-matching procedure.}
Ground-truth trajectories are map-matched to the road graph using nearest-node snapping.
For non-adjacent consecutive nodes (GPS jumps), we insert the shortest-path bridge between them.
Routes with excessive bridging (indicating poor GPS quality) are filtered.

\section{Waypoint Extraction Details}
\label{app:waypoints}

We extract waypoints from ground-truth graph paths using a modified Ramer--Douglas--Peucker (RDP) algorithm.
Given a path $(v_0, \dots, v_L)$, we iteratively select the node with maximum perpendicular distance from the current polyline approximation.
We continue until $M$ waypoints are selected (including origin and destination).
In our experiments, we use $M=4$ (origin + 2 intermediate + destination).

\section{Implementation Details}
\label{app:impl}

\paragraph{Waypoint AR model.}
The waypoint predictor is a 4-layer transformer encoder with $d_{\mathrm{model}}=128$ and 4 attention heads.
Input features include:
\begin{itemize}
  \item Waypoint bin embeddings (learned, $32 \times 32 = 1024$ classes)
  \item Destination bin embedding
  \item Hour-of-day (sinusoidal encoding, 24 bins)
  \item Day-of-week (one-hot, 7 classes)
  \item Road-tier context (aggregated from path statistics)
\end{itemize}
We train with cross-entropy loss using AdamW ($\text{lr}=10^{-3}$, weight decay $10^{-4}$) for 50 epochs with batch size 256.

\paragraph{A* execution.}
We use edge-length as the cost function and Euclidean distance to destination as the heuristic.
Maximum expansion budget is set to 10,000 nodes per waypoint segment.

\section{Additional Results}
\label{app:results}

\paragraph{Per-city breakdown.}
\Cref{tab:per-city} shows corridor generation results separately for Detroit and Columbus.
Columbus exhibits higher success rates due to denser road networks and shorter average path lengths.

\begin{table}[h]
  \caption{Per-city corridor generation results.}
  \label{tab:per-city}
  \centering
  \small
  \begin{tabular}{llcc}
    \toprule
    City & Method & Success & $J_{\mathrm{best}}$ \\
    \midrule
    Detroit & Node AR & 0.312 & 0.189 \\
    Detroit & Waypoint AR + A* & 1.000 & 0.251 \\
    \midrule
    Columbus & Node AR & 0.404 & 0.213 \\
    Columbus & Waypoint AR + A* & 1.000 & 0.281 \\
    \bottomrule
  \end{tabular}
\end{table}

\paragraph{Effect of bin granularity.}
Finer bins ($64 \times 64$) improve corridor match slightly ($J_{\mathrm{best}}$ mean $+0.02$) but reduce waypoint prediction accuracy due to sparser supervision per bin.
Coarser bins ($16 \times 16$) degrade corridor resolution.
The $32 \times 32$ setting balances these trade-offs for our dataset size.
